\section{An exact eighteen-qubit balanced-purity minimum}
\label{sec:eighteen}

\begin{theorem}[exact Rains dual and graph certificate; E18-1--E18-4]
\label{thm:eighteen}
For every normalised pure $18$-qubit state,
\[
 \boxed{\quad\pi_9(\psi)\ge\frac7{2288},\qquad
 \min_\psi\pi_9(\psi)=\frac7{2288}.\quad}
\]
Equality is attained by the graph state whose eighteen
adjacency rows (vertex $0$ is the lowest bit, in hexadecimal)
are
\[
\begin{gathered}
\texttt{36e72 0fb7d 1b80a 0cad6 25e0b 30583}\\
\texttt{1c10b 2b228 1a262 3999b 27031 3221f}\\
\texttt{2a696 1dd87 2a45b 273ce 22b65 1deb1}.
\end{gathered}
\]
Its entire graph-state Pauli weight distribution is
\[
 A_0=1,\quad
 (A_8,A_{10},A_{12},A_{14},A_{16},A_{18})
 =(2754,18360,77112,110160,50949,2808)
\]
and all other weights vanish. Its $9|9$ cut-rank
distribution is $7^{9180}9^{39440}$.
The graph is of Type II (all vertex degrees odd) and has
global spin-flip overlap $\tau_{18}=1$.
\end{theorem}
\begin{proof}
E18-1 verifies the rational dual identity for the
18-qubit Rains LP, with constant $7/2288$.
Specifically, using $a_s=a_s(A)$ and shadow coefficients
$S_i(A)$, it verifies coefficientwise
\[
 a_9-\tfrac7{2288}A_0
 =\sum_{s=0}^{8}e_s(a_s-a_{18-s})
 +\tfrac1{1166880}S_4
 +\tfrac1{1633632}S_6
 +\tfrac1{2498496}S_{11}
 +\sum_{j=0}^{18}w_jA_j
\]
with $w_j\ge0$ and the rational multipliers
\[
 (e_0,\ldots,e_8)=
 \left(\tfrac5{4004},-\tfrac{291}{2002},\tfrac{81}{364},
 \tfrac{278}{1001},-\tfrac{333}{1001},0,
 -\tfrac{101}{143},-\tfrac{54}{143},\tfrac{51}{26}\right).
\]
For a pure state the complementarity terms vanish and
every remaining term is nonnegative.

E18-2 checks the displayed binary graph is simple,
symmetric and has odd degrees; exhaustive enumeration of
its $2^{18}$ stabilizers gives the displayed weight
distribution, corresponding to an extremal Type II
self-dual additive $GF(4)$ code of minimum weight eight.
E18-3 computes every one of the $\binom{18}{9}=48620$
binary cut ranks, obtaining $9180$ cuts of rank $7$
and $39440$ of rank $9$. Thus
\[
 \pi_9=
 \frac{9180/128+39440/512}{48620}=\frac7{2288}.
\]
E18-4 independently checks the same purity from the
Pauli weight formula. The matching lower and upper
bounds prove the theorem.
\end{proof}

\begin{remark}[relationship to the known quaternary $S_{18}$ code]
An extremal Hermitian self-dual quaternary
$[18,9,8]$ code is classical and well documented
(see Cheng--Sloane and Harada--Munemasa's
\href{https://www.math.is.tohoku.ac.jp/~mharada/research/codes/hsd4.htm}
{code database}).
The graph above is a local-Clifford graph form of a
generator in that database. E18-1--E18-4 make the
purity/cut-rank attainment claim entirely reproducible
without downloading the external generator.
We do not claim novelty for the underlying code or
the existence of an extremal $[18,9,8]$ code.
The equality between its purity and the Rains-LP
lower bound is the conclusion here.
\end{remark}

\begin{corollary}[further sharp time-reversal-symmetric minima;
TY-1--TY-2]\label{cor:typeII-sixteen-twenty}
Inside the full spin-flip-invariant marginal class,
\[
 \boxed{\quad
 \min_{\tau_{16}=1}\pi_8=\frac1{143},\quad
 \min_{\tau_{18}=1}\pi_9=\frac7{2288},\quad
 \min_{\tau_{20}=1}\pi_{10}=\frac{131}{77792}.
 \quad}
\]
For $n=16$ and $20$ the respective attaining graph
states are the circulants $C_{16}(2,3,4,8)$ and
$C_{20}(3,4,5,10)$. Their balanced-cut rank
distributions are
\[
\begin{array}{c|l|c}
n & \text{ranks} & \pi_{n/2}\\\hline
16&6^{1096}7^{6882}8^{4892}&1/143\\
20&8^{12312}9^{96900}10^{75544}&131/77792.
\end{array}
\]
In particular the positive spin-flip coefficients
$c_{16}=1/1716$ and $c_{20}=103/1478048$
in Theorem~\ref{thm:purity-spinflip-tradeoff} are
\emph{sharp with their stated constant terms}.
The $n=20$ constant $1193/739024$ remains only a
global Rains lower bound, not a proven global minimum.
\end{corollary}
\begin{proof}
Theorem~\ref{thm:purity-spinflip-tradeoff} at
$\tau_n=1$ gives the lower values $1/143$
and $131/77792$ for $n=16,20$. TY-1--TY-2
check all $2^n$ graph stabilizers and all
balanced-cut ranks, with exact fraction arithmetic.
Each circulant has odd degree at every vertex, hence
$\tau_n=1$. Theorem~\ref{thm:eighteen} handles
$n=18$ and also exhibits a $\tau_{18}=1$ minimiser.
A stronger coefficient at $16$ or $20$ would fail
on the respective circulant witness, proving its
sharpness.
\end{proof}
